% part: intuitionistic-logic
% chapter: soundness-completeness
% section: canonical-model

\documentclass[../../../include/open-logic-section]{subfiles}

\begin{document}

\olfileid[tr]{int}{sc}{mod}

\olsection{Kanonik Model}

Modelimizin dünyaları, doğal sayıların sonlu dizileri~$\sigma$, yani
$\sigma\in\Nat^*$ olacaktır. $\Nat^*$'ın tümevarımla şöyle tanımlandığına
dikkat edin:
\begin{enumerate}
\item $\emptyseq\in\Nat^*$.
\item $\sigma\in\Nat^*$ ve $n\in\Nat$ ise $\sigma.n\in\Nat^*$ (burada
  $\sigma.n$, $\sigma\concat\tuple{n}$'dir; $\sigma\concat\sigma'$ ise
  $\sigma$ ile $\sigma'$ dizilerinin art arda eklenmesidir).
\item Bunların dışında hiçbir şey $\Nat^*$ içinde değildir.
\end{enumerate}
Dolayısıyla tümevarımlı tanımlar vermek için $\Nat^*$'ı kullanabiliriz.

$\tuple{!B_1,!C_1}$, $\tuple{!B_2,!C_2}$, \dots, bütün !!{formula}
çiftlerinin bir numaralandırması olsun. Bir $\Delta$ !!{formula} kümesi
verildiğinde $\Delta(\sigma)$'yı tümevarımla şöyle tanımlayın:
\begin{enumerate}
\item $\Delta(\emptyseq)=\Delta$
\item $\Delta(\sigma.n)={}$
  \[
  \begin{cases}
    (\Delta(\sigma) \cup \{!B_n\})^* &
    \text{$\Delta(\sigma)\cup\{!B_n\}\Proves/ !C_n$ ise} \\
    \Delta(\sigma) & \text{aksi hâlde.}
  \end{cases}
  \]
\end{enumerate}
Burada $(\Delta(\sigma)\cup\{!B_n\})^*$ ile,
\olref[lin]{lem:lindenbaum}'ın $\Delta(\sigma)\cup\{!B_n\}$ kümesine
ve $!C_n$ !!{formula}{}üne uygulanmasıyla varlığı elde edilen asal
!!{formula} kümesini kastediyoruz. Bu tanım gereğince
$\Delta(\sigma)\cup\{!B_n\}\Proves/ !C_n$ ise
$\Delta(\sigma.n)\Proves !B_n$ ve $\Delta(\sigma.n)\Proves/ !C_n$.
Ayrıca her~$n$ için $\Delta(\sigma)\subseteq\Delta(\sigma.n)$.
$\Delta$ asalsa her~$\sigma$ için $\Delta(\sigma)$ da asaldır.

\begin{defn}\ollabel{defn:canonical-model}
  $\Delta$ asal olsun. $\Delta$ için \emph{kanonik model}
  $\mModel{M(\Delta)}$ şöyle tanımlanır:
  \begin{enumerate}
  \item $W=\Nat^*$, doğal sayıların sonlu dizileri kümesidir.
  \item $R$, $R\sigma\sigma'$ ancak ve ancak $\sigma$, $\sigma'$'nün bir
    başlangıç parçası olduğunda (yani bazı $\sigma''$ dizisi için
    $\sigma'=\sigma\concat\sigma''$ olduğunda) geçerli olan kısmi sıradır.
  \item $V(p)=\Setabs{\sigma}{p\in\Delta(\sigma)}$.
  \end{enumerate}
\end{defn}

$R$'nin gerçekten bir kısmi sıra olduğu kolayca doğrulanır. $V$ üzerindeki
  monotonluk koşulu da sağlanır. $\Delta(\sigma)\subseteq\Delta(\sigma.n)$
  olduğundan, $\sigma'$yi $\sigma$'dan ayıran sonlu uzantının uzunluğu
  üzerinde tümevarımla, $R\sigma\sigma'$ olduğunda
  $\Delta(\sigma)\subseteq\Delta(\sigma')$ elde ederiz.

\end{document}
